\documentclass[11pt]{article}

\usepackage[margin=1in]{geometry}
\usepackage{amsmath,amssymb,mathtools}
\usepackage{booktabs}
\usepackage{array}
\usepackage{hyperref}

\setlength{\parindent}{0pt}
\setlength{\parskip}{0.55em}
\renewcommand{\arraystretch}{1.18}

\newcommand{\dd}{\,\mathrm{d}}
\newcommand{\ii}{\mathrm{i}}
\newcommand{\eoc}{\mathrm{EOC}}
\newcommand{\Oh}{\mathcal{O}}
\newcommand{\norm}[1]{\left\lVert #1\right\rVert}

\title{Mathematical Derivation Handout\\
Finite Difference Methods for the One-Way Wave Equation}
\author{Rishabh Kumar}
\date{\today}

\begin{document}
\maketitle

\section{Model Problem and Notation}

We consider
\begin{equation}
    u_t+a u_x=0,\qquad a>0,
\end{equation}
with the assignment value $a=1$.  The numerical grid is
\begin{equation}
    x_m=-2+mh,\qquad t_n=nk,\qquad v_m^n\approx u(x_m,t_n),
\end{equation}
and the Courant number is
\begin{equation}
    \lambda=\frac{k}{h},\qquad \nu=a\lambda=\frac{ak}{h}.
\end{equation}
For the present problem $a=1$, hence $\nu=\lambda$.

The exact solution follows from characteristics:
\begin{equation}
    \frac{\dd x}{\dd t}=a,\qquad \frac{\dd u}{\dd t}=0.
\end{equation}
Thus $x-at$ is constant and
\begin{equation}
    u(x,t)=
    \begin{cases}
        u_0(x-at), & x-at\ge -2,\\
        0, & x-at<-2,
    \end{cases}
\end{equation}
where
\begin{equation}
    u_0(x)=
    \begin{cases}
        (1-|x|)^2, & |x|\le 1,\\
        0, & |x|>1.
    \end{cases}
\end{equation}

\section{Discrete Norms and Experimental Order}

Let $T=t_N$ be the final time of a computation, and define the grid error by
\begin{equation}
    e_m^N=v_m^N-u(x_m,T).
\end{equation}
The three error measures used in the numerical tables and plots are
\begin{align}
    E_1
    &=
    h\sum_m |e_m^N|,
    \\
    E_2
    &=
    \left(h\sum_m |e_m^N|^2\right)^{1/2},
    \\
    E_\infty
    &=
    \max_m |e_m^N|.
\end{align}
Here $E_1$ and $E_2$ approximate integral norms:
\begin{equation}
    E_1\approx \norm{e(\cdot,T)}_{L^1},
    \qquad
    E_2\approx \norm{e(\cdot,T)}_{L^2}.
\end{equation}
The maximum norm $E_\infty$ is a pointwise norm.  It records the largest error at any
single grid point, and is therefore more sensitive to localized oscillations,
smearing, or errors near nonsmooth features.

For two successive meshes $h$ and $h/2$, the experimental order of convergence is
\begin{equation}
    \eoc
    =
    \frac{\log(E_h/E_{h/2})}{\log 2}.
\end{equation}
If $E_h\approx C h^p$, then $\eoc\approx p$.

For the $\lambda$-sweep stability plots, we also use the final-time amplitude
\begin{equation}
    \norm{v^N}_\infty=\max_m |v_m^N|.
\end{equation}
Since the exact transported pulse satisfies $0\le u(x,t)\le1$, very large values of
$\norm{v^N}_\infty$ are numerical evidence of instability rather than physical growth.

\section{Definition of Local Truncation Error}

The local truncation error is the residual obtained by inserting the exact solution
into the finite difference scheme.  If a method can be written in divided residual
form as
\begin{equation}
    L_{h,k}v_m^n=0,
\end{equation}
then the local truncation error is
\begin{equation}
    \tau_m^n=(L_{h,k}u)_m^n,
\end{equation}
where $u$ solves the PDE exactly.  Thus $\tau_m^n$ measures how well the finite
difference equation represents the differential equation at one grid point.

If
\begin{equation}
    \tau_m^n=\Oh(k^p+h^q),
\end{equation}
the scheme is locally $p$th order in time and $q$th order in space.  For the
advection problem, convergence studies usually keep
\begin{equation}
    \lambda=\frac{k}{h}
\end{equation}
fixed, so $k=\lambda h=\Oh(h)$.  Under this hyperbolic scaling,
\begin{equation}
    \Oh(k+h)=\Oh(h),\qquad
    \Oh(k+h^2)=\Oh(h),\qquad
    \Oh(k^2+h^2)=\Oh(h^2).
\end{equation}

Local truncation error is a consistency measure, not the final numerical error by
itself.  A small local residual must be combined with stability to obtain convergence.
This is why FTCS can be locally consistent but still unstable for pure advection.

\section{Taylor Expansions Used for Local Truncation Error}

The local truncation error is obtained by substituting the exact solution into a finite difference scheme.  The following expansions are repeatedly used:
\begin{align}
    u(x_m,t_n+k)
    &=u+k u_t+\frac{k^2}{2}u_{tt}+\frac{k^3}{6}u_{ttt}+\Oh(k^4),\\
    u(x_m,t_n-k)
    &=u-k u_t+\frac{k^2}{2}u_{tt}-\frac{k^3}{6}u_{ttt}+\Oh(k^4),\\
    u(x_m+h,t_n)
    &=u+h u_x+\frac{h^2}{2}u_{xx}+\frac{h^3}{6}u_{xxx}+\Oh(h^4),\\
    u(x_m-h,t_n)
    &=u-h u_x+\frac{h^2}{2}u_{xx}-\frac{h^3}{6}u_{xxx}+\Oh(h^4).
\end{align}
All derivatives above are evaluated at $(x_m,t_n)$ unless otherwise stated.

The PDE gives useful derivative identities:
\begin{equation}
    u_t=-a u_x,\qquad
    u_{tt}=a^2u_{xx},\qquad
    u_{ttt}=-a^3u_{xxx}.
\end{equation}

\section{Local Truncation Error: Explicit Schemes}

\subsection{FTFS}

The FTFS scheme is
\begin{equation}
    \frac{v_m^{n+1}-v_m^n}{k}
    +a\frac{v_{m+1}^n-v_m^n}{h}=0.
\end{equation}
Substituting the exact solution gives
\begin{align}
    \tau_m^n
    &=
    \left(u_t+\frac{k}{2}u_{tt}+\Oh(k^2)\right)
    +a\left(u_x+\frac{h}{2}u_{xx}+\Oh(h^2)\right)\\
    &=u_t+a u_x+\frac{k}{2}u_{tt}+\frac{ah}{2}u_{xx}
    +\Oh(k^2+h^2).
\end{align}
Since $u_t+a u_x=0$,
\begin{equation}
    \tau_m^n
    =
    \frac{k}{2}u_{tt}+\frac{ah}{2}u_{xx}
    +\Oh(k^2+h^2)
    =\Oh(k+h).
\end{equation}
Thus FTFS is formally first-order accurate when $k=\Oh(h)$.

\subsection{FTBS}

The FTBS scheme is
\begin{equation}
    \frac{v_m^{n+1}-v_m^n}{k}
    +a\frac{v_m^n-v_{m-1}^n}{h}=0.
\end{equation}
The backward spatial difference satisfies
\begin{equation}
    \frac{u(x_m,t_n)-u(x_m-h,t_n)}{h}
    =
    u_x-\frac{h}{2}u_{xx}+\Oh(h^2).
\end{equation}
Therefore
\begin{align}
    \tau_m^n
    &=
    u_t+a u_x+\frac{k}{2}u_{tt}
    -\frac{ah}{2}u_{xx}
    +\Oh(k^2+h^2)\\
    &=
    \frac{k}{2}u_{tt}
    -\frac{ah}{2}u_{xx}
    +\Oh(k^2+h^2)
    =\Oh(k+h).
\end{align}
FTBS is also first-order accurate.  For $a>0$ it is the correct upwind direction.

\subsection{FTCS}

The FTCS scheme is
\begin{equation}
    \frac{v_m^{n+1}-v_m^n}{k}
    +a\frac{v_{m+1}^n-v_{m-1}^n}{2h}=0.
\end{equation}
The centered spatial difference is
\begin{equation}
    \frac{u(x_m+h,t_n)-u(x_m-h,t_n)}{2h}
    =
    u_x+\frac{h^2}{6}u_{xxx}+\Oh(h^4).
\end{equation}
Hence
\begin{equation}
    \tau_m^n
    =
    \frac{k}{2}u_{tt}
    +\frac{a h^2}{6}u_{xxx}
    +\Oh(k^2+h^4)
    =
    \Oh(k+h^2).
\end{equation}
For the usual hyperbolic scaling $k=\Oh(h)$, FTCS is first-order consistent.

\subsection{CTCS or Leapfrog}

The CTCS scheme is
\begin{equation}
    \frac{v_m^{n+1}-v_m^{n-1}}{2k}
    +a\frac{v_{m+1}^n-v_{m-1}^n}{2h}=0.
\end{equation}
The centered time difference gives
\begin{equation}
    \frac{u(x_m,t_n+k)-u(x_m,t_n-k)}{2k}
    =
    u_t+\frac{k^2}{6}u_{ttt}+\Oh(k^4).
\end{equation}
Therefore
\begin{equation}
    \tau_m^n
    =
    \frac{k^2}{6}u_{ttt}
    +\frac{a h^2}{6}u_{xxx}
    +\Oh(k^4+h^4)
    =
    \Oh(k^2+h^2).
\end{equation}
CTCS is second-order accurate for smooth solutions.

\subsection{Lax--Friedrichs}

The Lax--Friedrichs scheme can be written as
\begin{equation}
    \frac{v_m^{n+1}-\frac12(v_{m+1}^n+v_{m-1}^n)}{k}
    +a\frac{v_{m+1}^n-v_{m-1}^n}{2h}=0.
\end{equation}
The spatial averaging term satisfies
\begin{equation}
    \frac12\left(u(x_m+h,t_n)+u(x_m-h,t_n)\right)
    =
    u+\frac{h^2}{2}u_{xx}+\Oh(h^4).
\end{equation}
Thus
\begin{align}
    \frac{u(x_m,t_n+k)-\frac12(u(x_m+h,t_n)+u(x_m-h,t_n))}{k}
    &=
    u_t+\frac{k}{2}u_{tt}
    -\frac{h^2}{2k}u_{xx}
    +\Oh(k^2+h^4/k).
\end{align}
Combining this with the centered spatial derivative,
\begin{equation}
    \tau_m^n
    =
    \frac{k}{2}u_{tt}
    -\frac{h^2}{2k}u_{xx}
    +\frac{a h^2}{6}u_{xxx}
    +\Oh(k^2+h^4/k+h^4).
\end{equation}
Therefore
\begin{equation}
    \tau_m^n=\Oh\left(k+\frac{h^2}{k}\right).
\end{equation}
Under $k=\Oh(h)$, Lax--Friedrichs is first-order accurate.

\subsection{Lax--Wendroff}

The Lax--Wendroff method is
\begin{equation}
    v_m^{n+1}
    =
    v_m^n
    -\frac{\nu}{2}(v_{m+1}^n-v_{m-1}^n)
    +\frac{\nu^2}{2}(v_{m+1}^n-2v_m^n+v_{m-1}^n).
\end{equation}
Equivalently,
\begin{equation}
    \frac{v_m^{n+1}-v_m^n}{k}
    +a\frac{v_{m+1}^n-v_{m-1}^n}{2h}
    -\frac{a^2k}{2}\frac{v_{m+1}^n-2v_m^n+v_{m-1}^n}{h^2}=0.
\end{equation}
Using
\begin{equation}
    \frac{u(x_m+h)-2u(x_m)+u(x_m-h)}{h^2}
    =
    u_{xx}+\frac{h^2}{12}u_{xxxx}+\Oh(h^4),
\end{equation}
we obtain
\begin{align}
    \tau_m^n
    &=
    u_t+a u_x
    +\frac{k}{2}u_{tt}
    -\frac{a^2k}{2}u_{xx}
    +\frac{k^2}{6}u_{ttt}
    +\frac{a h^2}{6}u_{xxx}
    +\Oh(k^3+kh^2+h^4).
\end{align}
The terms $u_t+a u_x$ vanish by the PDE, and
$u_{tt}=a^2u_{xx}$ cancels the order-$k$ term.  Hence
\begin{equation}
    \tau_m^n=\Oh(k^2+h^2).
\end{equation}

\subsection{Beam--Warming}

For $a>0$, the Beam--Warming method uses a second-order backward stencil:
\begin{equation}
    v_m^{n+1}
    =
    v_m^n
    -\frac{\nu}{2}(3v_m^n-4v_{m-1}^n+v_{m-2}^n)
    +\frac{\nu^2}{2}(v_m^n-2v_{m-1}^n+v_{m-2}^n).
\end{equation}
The needed backward approximations are
\begin{align}
    \frac{3u_m-4u_{m-1}+u_{m-2}}{2h}
    &=
    u_x-\frac{h^2}{3}u_{xxx}+\Oh(h^3),\\
    \frac{u_m-2u_{m-1}+u_{m-2}}{h^2}
    &=
    u_{xx}-h u_{xxx}+\Oh(h^2).
\end{align}
Substitution gives
\begin{align}
    \tau_m^n
    &=
    u_t+a u_x
    +\frac{k}{2}u_{tt}
    -\frac{a^2k}{2}u_{xx}
    +\Oh(k^2+h^2+kh).
\end{align}
Again $u_t+a u_x=0$ and $u_{tt}=a^2u_{xx}$, so
\begin{equation}
    \tau_m^n=\Oh(k^2+h^2)
\end{equation}
under the hyperbolic scaling $k=\Oh(h)$.

\section{Local Truncation Error: Implicit Schemes}

\subsection{BTFS and BTBS}

BTFS is
\begin{equation}
    \frac{v_m^{n+1}-v_m^n}{k}
    +a\frac{v_{m+1}^{n+1}-v_m^{n+1}}{h}=0,
\end{equation}
while BTBS is
\begin{equation}
    \frac{v_m^{n+1}-v_m^n}{k}
    +a\frac{v_m^{n+1}-v_{m-1}^{n+1}}{h}=0.
\end{equation}
Expanding around $(x_m,t_{n+1})$,
\begin{equation}
    \frac{u(x_m,t_{n+1})-u(x_m,t_n)}{k}
    =
    u_t-\frac{k}{2}u_{tt}+\Oh(k^2).
\end{equation}
The forward and backward spatial differences are first-order accurate.  Therefore
\begin{equation}
    \tau_m^{n+1}=\Oh(k+h)
\end{equation}
for both BTFS and BTBS.

\subsection{BTCS}

BTCS uses a backward time difference and a centered spatial difference at the new time:
\begin{equation}
    \frac{v_m^{n+1}-v_m^n}{k}
    +a\frac{v_{m+1}^{n+1}-v_{m-1}^{n+1}}{2h}=0.
\end{equation}
Thus
\begin{equation}
    \tau_m^{n+1}
    =
    -\frac{k}{2}u_{tt}
    +\frac{a h^2}{6}u_{xxx}
    +\Oh(k^2+h^4)
    =
    \Oh(k+h^2).
\end{equation}
With $k=\Oh(h)$, BTCS is first-order accurate.

\subsection{Crank--Nicolson}

The Crank--Nicolson method is
\begin{equation}
    \frac{v_m^{n+1}-v_m^n}{k}
    +\frac{a}{2}
    \left[
        \frac{v_{m+1}^{n+1}-v_{m-1}^{n+1}}{2h}
        +
        \frac{v_{m+1}^{n}-v_{m-1}^{n}}{2h}
    \right]=0.
\end{equation}
This is the trapezoidal rule in time applied to the centered spatial derivative.  Expanding about the midpoint $t_{n+1/2}$ gives second-order time accuracy, while the centered spatial derivative is second-order in $h$. Hence
\begin{equation}
    \tau_m^{n+1/2}=\Oh(k^2+h^2).
\end{equation}

\section{Von Neumann Stability: Explicit Schemes}

For stability, insert the Fourier mode
\begin{equation}
    v_m^n=G^n e^{\ii m\theta}.
\end{equation}
A scheme is stable in the von Neumann sense if $|G(\theta)|\le 1$ for every wave number $\theta$.

\subsection{FTFS}

FTFS gives
\begin{equation}
    G=1-\nu(e^{\ii\theta}-1)=1+\nu-\nu e^{\ii\theta}.
\end{equation}
Therefore
\begin{align}
    |G|^2
    &=
    (1+\nu-\nu\cos\theta)^2+\nu^2\sin^2\theta\\
    &=
    1+2\nu(1+\nu)(1-\cos\theta).
\end{align}
For $\nu>0$ and $\theta\ne 0$, $|G|>1$.  Thus FTFS is unstable for right-moving advection.

\subsection{FTBS}

FTBS gives
\begin{equation}
    G=1-\nu(1-e^{-\ii\theta})
    =
    1-\nu+\nu e^{-\ii\theta}.
\end{equation}
Hence
\begin{align}
    |G|^2
    &=
    (1-\nu+\nu\cos\theta)^2+\nu^2\sin^2\theta\\
    &=
    1-2\nu(1-\nu)(1-\cos\theta).
\end{align}
Thus $|G|\le 1$ for all $\theta$ precisely when
\begin{equation}
    0\le \nu\le 1.
\end{equation}

\subsection{FTCS}

FTCS gives
\begin{equation}
    G=1-\frac{\nu}{2}(e^{\ii\theta}-e^{-\ii\theta})
    =
    1-\ii\nu\sin\theta.
\end{equation}
Thus
\begin{equation}
    |G|^2=1+\nu^2\sin^2\theta.
\end{equation}
For any nonzero $\nu$ and any mode with $\sin\theta\ne0$, $|G|>1$.  FTCS is therefore unstable for pure advection.

\subsection{Lax--Friedrichs}

For Lax--Friedrichs,
\begin{equation}
    G=\frac12(e^{\ii\theta}+e^{-\ii\theta})
    -\frac{\nu}{2}(e^{\ii\theta}-e^{-\ii\theta})
    =
    \cos\theta-\ii\nu\sin\theta.
\end{equation}
Consequently
\begin{equation}
    |G|^2=\cos^2\theta+\nu^2\sin^2\theta
    =
    1-(1-\nu^2)\sin^2\theta.
\end{equation}
The condition $|G|\le1$ for all $\theta$ is
\begin{equation}
    |\nu|\le1.
\end{equation}

\subsection{Lax--Wendroff}

The Lax--Wendroff amplification factor is
\begin{equation}
    G
    =
    1-\ii\nu\sin\theta+\nu^2(\cos\theta-1).
\end{equation}
Let $s=\sin(\theta/2)$.  Then
\begin{equation}
    \cos\theta-1=-2s^2,\qquad
    \sin^2\theta=4s^2(1-s^2).
\end{equation}
Therefore
\begin{align}
    |G|^2
    &=
    (1-2\nu^2s^2)^2+4\nu^2s^2(1-s^2)\\
    &=
    1-4\nu^2(1-\nu^2)s^4.
\end{align}
Thus $|G|\le1$ for all $\theta$ if and only if
\begin{equation}
    |\nu|\le1.
\end{equation}

\subsection{Beam--Warming}

Let $q=e^{-\ii\theta}$.  Beam--Warming gives
\begin{equation}
    G=
    1-\frac{\nu}{2}(3-4q+q^2)
    +\frac{\nu^2}{2}(1-2q+q^2).
\end{equation}
Set $\delta=1-q$.  Since
\begin{equation}
    3-4q+q^2=2\delta+\delta^2,\qquad
    1-2q+q^2=\delta^2,
\end{equation}
we can write
\begin{equation}
    G=1-\nu\delta+\frac{\nu(\nu-1)}{2}\delta^2.
\end{equation}
Using
\begin{equation}
    \delta+\overline{\delta}=|\delta|^2=4\sin^2(\theta/2),
\end{equation}
one obtains after simplification
\begin{equation}
    |G|^2
    =
    1-4\nu(2-\nu)(1-\nu)^2\sin^4(\theta/2).
\end{equation}
Hence
\begin{equation}
    |G|\le1 \quad\text{for all }\theta
    \qquad\Longleftrightarrow\qquad
    0\le\nu\le2.
\end{equation}

\subsection{CTCS}

CTCS gives
\begin{equation}
    G-G^{-1}=-\nu(e^{\ii\theta}-e^{-\ii\theta})
    =
    -2\ii\nu\sin\theta.
\end{equation}
Multiplying by $G$,
\begin{equation}
    G^2+2\ii\nu\sin\theta\,G-1=0.
\end{equation}
The roots are
\begin{equation}
    G=-\ii\nu\sin\theta
    \pm\sqrt{1-\nu^2\sin^2\theta}.
\end{equation}
For the roots to remain on the unit circle for all $\theta$, the square root must be real for all $\theta$, which requires
\begin{equation}
    |\nu|\le1.
\end{equation}
This is the CFL condition for CTCS.  The method also supports a computational mode, which is why a stable one-step startup is needed.

\section{Von Neumann Stability: Implicit Schemes}

\subsection{BTFS}

BTFS gives
\begin{equation}
    G\left[1+\nu(e^{\ii\theta}-1)\right]=1,
\end{equation}
so
\begin{equation}
    G=\frac{1}{1-\nu+\nu e^{\ii\theta}}.
\end{equation}
The denominator has modulus squared
\begin{equation}
    |1-\nu+\nu e^{\ii\theta}|^2
    =
    1-2\nu(1-\nu)(1-\cos\theta).
\end{equation}
For $0<\nu<1$, this can be less than $1$, so $|G|>1$ for some modes.  For $\nu\ge1$, the denominator modulus is at least $1$.  Thus BTFS is not the natural implicit upwind method for $a>0$ when $0<\nu<1$.

\subsection{BTBS}

BTBS gives
\begin{equation}
    G\left[1+\nu-\nu e^{-\ii\theta}\right]=1,
\end{equation}
so
\begin{equation}
    G=\frac{1}{1+\nu-\nu e^{-\ii\theta}}.
\end{equation}
The denominator satisfies
\begin{equation}
    |1+\nu-\nu e^{-\ii\theta}|^2
    =
    1+2\nu(1+\nu)(1-\cos\theta)\ge1.
\end{equation}
Therefore
\begin{equation}
    |G|\le1\qquad \text{for all }\nu\ge0.
\end{equation}
BTBS is unconditionally stable for right-moving advection, but it is diffusive.

\subsection{BTCS}

BTCS gives
\begin{equation}
    G\left[1+\frac{\nu}{2}(e^{\ii\theta}-e^{-\ii\theta})\right]=1.
\end{equation}
Thus
\begin{equation}
    G=\frac{1}{1+\ii\nu\sin\theta},
\end{equation}
and
\begin{equation}
    |G|^2=\frac{1}{1+\nu^2\sin^2\theta}\le1.
\end{equation}
So BTCS is unconditionally stable in the von Neumann sense.

\subsection{Crank--Nicolson}

Crank--Nicolson gives
\begin{equation}
    G\left[1+\frac{\ii\nu}{2}\sin\theta\right]
    =
    1-\frac{\ii\nu}{2}\sin\theta.
\end{equation}
Therefore
\begin{equation}
    G=
    \frac{1-\frac{\ii\nu}{2}\sin\theta}
         {1+\frac{\ii\nu}{2}\sin\theta}.
\end{equation}
The numerator and denominator are complex conjugates, so
\begin{equation}
    |G|=1.
\end{equation}
Crank--Nicolson is neutrally stable for every $\nu$.  This means it does not damp Fourier modes, which is useful for amplitude preservation but can also preserve dispersive oscillations near nonsmooth data.

\section{Summary of Orders and Stability}

\begin{table}[h]
\centering
\small
\begin{tabular}{@{}lll@{}}
\toprule
Scheme & Formal order for $k=\Oh(h)$ & Stability for $a>0$\\
\midrule
FTFS & first & unstable\\
FTBS & first & $0\le\nu\le1$\\
FTCS & first & unstable\\
CTCS & second & $|\nu|\le1$\\
Lax--Friedrichs & first & $|\nu|\le1$\\
Lax--Wendroff & second & $|\nu|\le1$\\
Beam--Warming & second & $0\le\nu\le2$\\
BTFS & first & stable for $\nu\ge1$ on a periodic grid\\
BTBS & first & all $\nu\ge0$\\
BTCS & first & all $\nu$\\
Crank--Nicolson & second & all $\nu$, neutrally stable\\
\bottomrule
\end{tabular}
\end{table}

The formal orders assume a smooth exact solution.  The assigned initial data has a corner at $x=0$, so the observed global convergence rates, especially in $E_\infty$, are lower than the smooth-solution stencil orders.  This explains why the numerical convergence plots show clear improvement under refinement but do not always display the full formal order.

\end{document}
